\begin{abstract}
When retrieved evidence becomes sufficient to answer a question, does a language model's internal state predict, before it generates a token, whether it will use that evidence? We run paired forward passes over \emph{evidence trajectories} that add distractors, individual reasoning hops, the decisive hop, redundant, falsified, partial or paraphrased support, or the answer string in an irrelevant passage, and read the residual stream at the answer-start position. Linear probes on the paired update $\Delta h = h(C_k)-h(C_{k-1})$ recover \emph{sufficiency} at $0.95$ AUROC on document-, answer- and relation-disjoint splits in three model families, but a fine-tuned text classifier recovers it as well. A nearly orthogonal direction predicts \emph{uptake}, whether decisive evidence will be integrated rather than ignored, at $0.78$ AUROC, where text classifiers stay at $0.52$--$0.63$: uptake is an internal signal. Causal tests separate readout from mechanism. Probe directions are inert under steering, projection-out and component patching; the full anchor residual is causally sufficient and its effect lies in the orthogonal complement; the only causally active direction, the sufficient-minus-insufficient mean difference, is an abstention gate that raises answering on distractors and false evidence as much as on sufficient evidence. These results replicate on HotpotQA, on RAGBench's naturally retrieved documents across five domains, on BM25- and E5-retrieved Wikipedia evidence, and on a synthetic benchmark, although Wikipedia-trained directions transfer only partially across domains. High-AUROC, transferable probes find the label, not the lever.
\end{abstract}
